\section{Six-point audit, claim boundary, and the next arithmetic doors}
\label{sec:tpc59-claims}

The critical cutoff changes the route in one precise way.  It replaces the
fixed-power displacement of the TPC--58 very-high face by an admissible
support choice at the same exponent as the inherited cofactor packet:
\begin{equation}
 y=C_{\rm vh}Q,
 \qquad
 R_{\rm safe}(y)\asymp R_{\rm max}(y)\asymp J.
 \label{eq:tpc59-final-critical-route}
\end{equation}
Together with the exact dyadic partition, this proves one support
feasibility statement and one zero fixed-power energy-localization cost:
\begin{equation}
 \text{the cutoff gate is open at }R\asymp J,
 \qquad
 \lambda_{\rm blk}=0
 \quad\text{after }\cD_H(y;\zeta)>0.
 \label{eq:tpc59-final-two-zero-costs}
\end{equation}
Neither statement is an occupancy or primality theorem.  Cutoff feasibility
is not itself an additive energy exponent.  The constants
\(\beta_{\cR}\), \(\rho_{\cR}(\zeta)\), and \(\tau_{\cR}\) remain the
first quantitative gates, and the undeleted singleton, native, and affine
gates remain downstream of them.  Restoring high Walsh coordinates outside
the represented terminal image is a further nonmonotone reassembly gate.

\subsection{Six-point physical audit}

We record the six checks that must be passed before an exact support or
linear-algebra statement is promoted to progress on the physical endpoint
route.  Passing these checks identifies the correct object; it does not
prove that any of the open lower constants is positive.

\begin{enumerate}[label=\textup{(A\arabic*)},leftmargin=3.2em]
 \item \emph{Literal physical coefficients.}
       The pre-terminal and terminal atoms are, respectively,
       \begin{equation}
        [G_X^L]_{m,r,j,\gamma}=c_{m,j,\gamma}\zeta_m,
        \qquad
        x_u=\mu(r_u)\Lambda_X(mj+h_0)c_{m,j,\gamma}\zeta_m.
        \label{eq:tpc59-audit-literal-coefficients}
       \end{equation}
       The coefficient \(c_{m,j,\gamma}\) retains the inherited source
       shape, M\"obius and prime-logarithm factors, masks, opposite-row
       factor, orbit phase, and four-factor profile.  The terminal multiplier
       is applied only at terminalization.  No abstract sign family,
       coefficient-blind support vector, or renormalized terminal atom is
       substituted for these physical coefficients.  The terminal/Walsh
       comparison uses the calibrated common gauge of TPC--58, so no second
       M\"obius factor is inserted.

 \item \emph{One fixed physical shift.}
       Every terminal coordinate obeys
       \begin{equation}
        mj+h_0=p_ur_u
        \label{eq:tpc59-audit-fixed-shift}
       \end{equation}
       for the one declared nonzero integer \(h_0\).  The cutoff theorem,
       largest-prime census, full-fiber collision relation, and endpoint
       audit are pointwise at this shift.  There is no average over shifts,
       intercepts, or frozen parameter families.

 \item \emph{Physical atomic normalization.}
       The parent energy \(\cE_{\rm par}\), compatible-window energy
       \(\cE_{\cR}\), terminal diagonal \(\cD_H\), and selected-block
       diagonal \(\cD_R\) all use raw counting measure on their declared
       atoms.  Repeated cofactor values in different coordinates remain
       distinct orthogonal atoms until the stated grouping.  In particular,
       the shaped weight \(\Lambda_X\) is not moved into a new input norm,
       and the atomic dyadic identity is not asserted for a coherent grouped
       square.

 \item \emph{Canonical representation and minimality scope.}
       The baseline retained-pair carrier \(\cI_X\), orbit cell, opposite-row
       cell, and cutoff are fixed independently of \(\zeta\).  On the
       very-high face, \(p_u>y>\sqrt{N_+}\) gives a unique largest-prime
       representation and hence a canonical cofactor \(r_u\).  The
       projection onto a fixed set \(\cR\) is therefore an actual coordinate
       projection, not a transport between cofactor labels.  By contrast,
       the maximal dyadic block \(k_*(\zeta)\) may depend on the distinguished
       physical vector.  Its pigeonhole estimate is a distinguished-vector
       statement and is not silently promoted to a coefficient-uniform
       fixed-block theorem.  The uniqueness assertion is canonicality of the
       very-high atomic representation on the declared one-side baseline; it
       is not a claim that the full coefficient parametrization is globally
       dimension-minimal.

 \item \emph{Actual active support.}
       Cutoff compatibility, cofactor occupancy, and terminal activation are
       kept separate.  Equation \eqref{eq:tpc59-final-critical-route} says
       only that exponent-\(133/400\) cofactors are allowed by the literal
       very-high inequalities.  The exact tests
       \begin{equation}
        \beta_{\cR}
        =\min_{m:B_mN(m)>0}\frac{N_{\cR}(m)}{N(m)},
        \qquad
        \tau_{\cR}
        =\min_{m:B_mN_{\cR}(m)>0}
          \frac{T_{\cR}(m)}{B_mN_{\cR}(m)}
        \label{eq:tpc59-audit-active-support-tests}
       \end{equation}
       may still vanish.  Moreover, the degree used in the singleton bridge
       is the degree in the complete active terminal vector.  Collision
       partners with \(p\le y\) are not deleted merely because the root lies
       on the very-high face.  The decomposition
       \eqref{eq:tpc59-base-high-partition} also keeps the missing high Walsh
       mass \(\cD_{\rm miss}\) separate from the represented terminal image.
       The TPC--58 high map used downstream is restricted to that represented
       image; putting \(\cD_{\rm miss}\) back requires its own reassembly
       estimate.

 \item \emph{Strict endpoint budget.}
       Cutoff feasibility is audited before the energy ledger.  Actual
       occupancy, terminal activation, singleton survival, the native Gram,
       represented affine completion, missing-high reassembly, and downstream
       boundary/tail transfers are separate ledger entries.  Only the
       conditional localization cost \(\lambda_{\rm blk}\) is proved zero
       here.  Every route to the endpoint must satisfy
       \begin{equation}
        \Lambda_{\rm tot}<\frac1{400}
        \label{eq:tpc59-audit-strict-budget}
       \end{equation}
       with \(\Lambda_{\rm tot}\) as in
       \eqref{eq:tpc59-complete-loss-ledger}.  Equality is not accepted
       without closed leading constants and quantified remainders.  The
       coefficient-blind capacity exponent in
       \eqref{eq:tpc59-compatible-capacity-budget-gap} exceeds the allowance;
       it closes that support-only certificate and is not entered as a
       proved physical singleton loss.
\end{enumerate}

\subsection{Claim table}

\begin{center}
\small
\begin{tabular}{>{\raggedright\arraybackslash}p{0.29\textwidth}
                >{\raggedright\arraybackslash}p{0.60\textwidth}}
\toprule
Statement & Status \\
\midrule
Critical admissible cutoff \(y=C_{\rm vh}Q\)
& Exact support theorem (L0), specialized to the literal fixed-\(h_0\)
  packet without changing its coefficients or normalization (L1). \\
Cutoff--cofactor reciprocity
& Exact relation
  \(R_{\rm safe}\asymp R_{\rm max}\asymp J/\omega_X\) and exact
  comparison of the critical, subpower, and fixed-power regimes (L0/L1).
  It is not an occupancy estimate. \\
Sharp cofactor-occupancy constant
& Exact rowwise minimax formula for \(\beta_{\cR}\), and exact
  distinguished ratio \(\rho_{\cR}(\zeta)\) (L0).  Positivity or a useful
  endpoint exponent for either quantity is not proved. \\
Terminal-activation constant
& Exact rowwise minimax formula for \(\tau_{\cR}\) and combined constant
  \(\chi_{\cR}\) (L0), using the literal terminal weights (L1).
  Positivity is not proved. \\
Largest-prime census
& Exact characterization of every very-high terminal atom by
  \(\Pplus(mj+h_0)\), including all retained-pair and terminal masks at one
  fixed shift (L1).  It is an identity, not a density theorem. \\
Dyadic localization
& Exact raw-energy partition and an \(X^{-o(1)}\) maximal-block loss after
  \(\cD_H>0\) (L0/L1).  Thus \(\lambda_{\rm blk}=0\) conditionally; no
  parent-to-high-face transfer is supplied. \\
Base/represented/missing decomposition
& Exact atomic decomposition including high Walsh mass absent from the
  retained terminal image (L0/L1).  The alternative favorable to a proof
  using the represented high face is not forced. \\
Core-to-high/full singleton bridge
& Exact lower floor for both the high-only residual square used downstream
  and the complete undeleted residual square, using roots singleton in the
  complete terminal fibers; every collision partner is localized to a
  constant-scale neighborhood (L1).  A positive weighted singleton fraction
  is not proved. \\
Critical endpoint support audit
& Exact proof that the cutoff feasibility gate is open at \(R\asymp J\);
  every fixed-power increase in the cutoff lowers the compatible cofactor
  exponent by the same power (L1 route correction).  Feasibility is not an
  additive energy loss. \\
Support-only singleton closures
& The inherited capacity bound already exceeds the \(1/400\) allowance,
  and the same-divisor and automatic-singleton thresholds lie outside the
  compatible carrier (L1 no-go for those specific certificates).  This is
  not a lower bound for the actual collision energy. \\
Native and affine transfers
& The exact downstream constants \(\alpha_H\) and \(\delta_{\rm aff}\)
  remain those of TPC--58 on the represented terminal/Walsh baseline.
  Their positivity and scale-uniform size are not consequences of critical
  cutoff tuning. \\
Missing-high reassembly
& Restoring \(\cW_H\setminus\Ran\iota\) after the represented affine
  transfer is a separate nonmonotone gate.  No positive reassembly constant
  or exponent is proved. \\
Endpoint conclusion
& Conditional exponent ledger only.  No positive fixed-\(h_0\) arithmetic
  lower square, parity improvement, prime-pair asymptotic, or twin-prime
  conclusion is proved (no L2 theorem). \\
\bottomrule
\end{tabular}
\end{center}

Here L0 denotes exact finite-dimensional identities, minimax laws, support
geometry, and sharp route-specific obstructions.  L1 denotes their exact
specialization to the inherited literal coefficients, raw atomic norm,
actual carrier, and one fixed nonzero shift.  An L2 result would require a
new quantitative arithmetic estimate for the distinguished physical vector
or the full canonical coefficient image, strong enough to survive every
arrow in \eqref{eq:tpc59-prospective-transfer-chain}.  No result in this
paper is classified as L2.

\subsection{The next seven doors and their kill tests}

The next work should follow the transfer order.  A downstream Gram
calculation cannot repair an empty compatible face, and a favorable atomic
block cannot be returned to the undeleted coherent vector by monotonicity.

\begin{enumerate}[label=\textbf{Door \Alph*.},leftmargin=3.8em]
 \item \emph{Critical-window occupancy: \(\beta_{\cR}\) or
       \(\rho_{\cR}(\zeta)\).}
       Fix a coefficient-independent critical cofactor set
       \(\cR_{\rm crit}\subset[c_-J,c_+J]\) allowed by
       \cref{thm:tpc59-critical-cutoff,thm:tpc59-cutoff-cofactor-reciprocity}.
       The uniform target is
       \begin{equation}
        \beta_{\cR_{\rm crit}}
        \ge X^{-\lambda_{\rm occ}+o(1)},
        \label{eq:tpc59-next-uniform-occupancy}
       \end{equation}
       while the weaker distinguished target is
       \begin{equation}
        \rho_{\cR_{\rm crit}}(\zeta)
        \ge X^{-\lambda_{\rm occ}+o(1)}.
        \label{eq:tpc59-next-distinguished-occupancy}
       \end{equation}
       These require information about the cofactors already present in the
       literal retained-pair carrier; there is no relocation map to analyze.
       The exact uniform kill test is an active row with
       \(N_{\cR_{\rm crit}}(m)=0\), which forces \(\beta=0\).  A
       distinguished route survives that failure only if its physical
       coefficient vanishes on every offending row and its own ratio remains
       quantitatively positive.  The adaptive block \(k_*(\zeta)\) may be
       used only after this distinction is stated.

 \item \emph{Terminal activation: \(\tau_{\cR}\).}
       On a surviving compatible set, prove
       \begin{equation}
        \tau_{\cR}
        \ge X^{-\lambda_{\rm term}+o(1)},
        \label{eq:tpc59-next-terminal-activation}
       \end{equation}
       or the explicitly labeled distinguished-vector analogue comparing
       \(\cD_H(y;\zeta)\) with \(\cE_{\cR}(\zeta)\).  The largest-prime
       census reduces the problem to the fixed-shift values
       \(\Pplus(mj+h_0)\), but does not estimate their density.  Every
       squarefree, coprimality, support, and shaped-terminal factor must stay
       in \(T_{\cR}(m)\).  An active compatible row with
       \(T_{\cR}(m)=0\) kills the coefficient-uniform route.  If the best
       activation loss consumes the remaining endpoint slack, the route
       stops before dyadic or Gram analysis.

 \item \emph{Weighted full-fiber singleton collision.}
       After \(\cD_H>0\), select the maximal atomic block and prove the
       literal one-sided estimate
       \begin{equation}
        \cI_{R_*}^{\rm full}(y;\zeta)
        \le(1-\delta_{R_*})\cD_{R_*}(y;\zeta),
        \qquad
        \delta_{R_*}\ge X^{-\lambda_{\rm sing}+o(1)}.
        \label{eq:tpc59-next-weighted-singleton}
       \end{equation}
       The degree is computed in the complete terminal fiber, and every
       partner in the enlarged block \([R_*/2,4R_*)\) is retained whether
       or not its prime exceeds \(y\).  Unweighted degree bounds, the
       coefficient-blind determinant capacity, and singleton counts inside
       the selected face alone do not prove this inequality.  If
       \(\delta_{R_*}=0\), or if every valid lower exponent already exceeds
       the remaining budget, the present singleton route closes.

 \item \emph{Native residual-label erasure: \(\alpha_H\).}
       Only after the undeleted residual energy has been obtained should one
       form the nonzero actual image \(E_H=\Ran T_H\) and seek
       \begin{equation}
        \alpha_H\ge X^{-\lambda_H+o(1)}.
        \label{eq:tpc59-next-native-Gram}
       \end{equation}
       The exact uniform kill test remains
       \begin{equation}
        E_H\cap\Ker C_{\mu,H}\ne\{0\}.
        \label{eq:tpc59-next-native-kernel-kill}
       \end{equation}
       A positive singleton fraction controls the first grouping
       \(p\mapsto(i,s)\); it does not control the second grouping
       \((i,s)\mapsto i\).  A finite-scale numerical eigenvalue is scouting
       evidence, not a scale-uniform lower certificate.

 \item \emph{Base--high affine transversality:
       \(\delta_{\rm aff}\).}
       Conditional on the native gate, form the nonzero represented feasible
       pair image, retain the joint dependence of the base and represented
       high columns on the same coefficient vector, and prove
       \begin{equation}
        \delta_{\rm aff}
        \ge X^{-\lambda_{\rm aff}+o(1)}.
        \label{eq:tpc59-next-affine-angle}
       \end{equation}
       The exact uniform kill test is a nonzero intersection of the actual
       feasible pair image with the anti-diagonal,
       \begin{equation}
        \cF_{\rm aff}\cap\{(g,-g):g\in\cH_{\rm out}\}
        \ne\{0\}.
        \label{eq:tpc59-next-affine-kernel-kill}
       \end{equation}
       Treating the base and high vectors as independent, or replacing the
       physical coefficient vector by a favorable adaptive pair, changes
       the problem.  If the uniform image meets the anti-diagonal, only an
       explicitly distinguished-vector route may remain.

 \item \emph{Missing-high/off-carrier reassembly.}
       Let \(g_{\rm rep}(\zeta)\) be the represented affine output after
       Door E and let \(h_{\rm miss}(\zeta)\) be the native contribution of
       the high Walsh coordinates in
       \(\cW_H(y)\setminus\Ran\iota\).  Both are driven by the same row
       coefficient vector.  The required lower transfer has the form
       \begin{equation}
        \|g_{\rm rep}(\zeta)+h_{\rm miss}(\zeta)\|_2^2
        \ge
        X^{-\lambda_{\rm reasm}+o(1)}
        \|g_{\rm rep}(\zeta)\|_2^2.
        \label{eq:tpc59-next-missing-high-reassembly}
       \end{equation}
       Its coefficient-uniform zero test is a nonzero actual pair
       \((g_{\rm rep},h_{\rm miss})=(g,-g)\).  The scalar
       \(\cD_{\rm miss}\) records omitted atomic mass but cannot rule out
       such vector cancellation.  If the anti-diagonal is met, a
       distinguished-vector estimate or a different carrier is required;
       simply restoring the omitted coordinates is not norm-monotone.

 \item \emph{Strict endpoint budget and route decision.}
       Insert every proved exponent into
       \begin{equation}
        \begin{aligned}
        \Lambda_{\rm tot}={}&
        \lambda_{\rm sel}+\lambda_{\rm occ}
        +\lambda_{\rm blk}
        +\lambda_{\rm term}+\lambda_{\rm sing}\\
        &+\lambda_H+\lambda_{\rm aff}
        +\lambda_{\rm reasm}
        +\lambda_{\rm bdry}+\lambda_{\rm tail},
        \end{aligned}
        \label{eq:tpc59-next-complete-budget}
       \end{equation}
       retaining \(\lambda_{\rm blk}=0\) only with the distinguished-vector
       scope in \eqref{eq:tpc59-final-two-zero-costs}; cutoff feasibility is
       checked before this sum.  Continue the endpoint
       route only if \(\Lambda_{\rm tot}<1/400\).  At equality, stop unless
       leading constants and all remainders are closed.  If a kernel or
       zero-occupancy test fires, or a necessary loss alone exhausts the
       slack, record the resulting route-specific obstruction rather than
       replacing the physical vector, shift, normalization, or carrier.
\end{enumerate}

\paragraph{Conclusion.}
The critical-cutoff analysis is a genuine L1 route correction.  It proves
that the earlier positive-power cofactor displacement was caused by a
nonoptimal cutoff schedule, splits geometry from actual occupancy, and
shows that dyadic localization costs no positive power after the very-high
face has been activated.  The surviving route is narrower and auditable:
 critical-window occupancy, terminal activation, weighted singleton
 survival, native erasure, represented affine transversality, missing-high
 reassembly, and the strict endpoint ledger must be passed in that order.
 None is bypassed here.  In
particular, this paper contains no L2 parity improvement, prime-pair
asymptotic, or twin-prime implication.
